\chapter{Introduction}
\label{intro}
\section{Topic}
\label{topic}
This project is part of the course Speech Recognition (ELEC-E5150) of the year 2026 at Aalto University and was chosen based on the competition description provided by the lecturer.

The aim of this project is the development of a speech recognition system for a ghibberish version of Esperanto based on the provided dataset.
Throughout the remainings of this project we will develop an automated system to transcribe Esperanto to its textual representation.
The used dataset contains artificially produced speech samples without any inherent meaning and is part of the publicly available Common Voice Corpus. 

Esperanto is a planned language that was invented by the polish doctor Ludwik Lejzer Zamenhof at the end of the 19th century to serve as a simple unified language for international communication.
In recent years it became obvious that Esperanto failed its purpose to serve as a \textit{world language} due to the rise and acceptance of english for nearly all international interactions.
However, there are still Esperanto speakers on every continent nowadays.
To enable as many people as possible to learn it easily, it was especially designed to have only very few grammatical rules without any exceptions and relies on a strong and fixed coupling between the written and spoken langage, what could make it better applicable for the aimed task.

\section{Problem analysis}
\label{problem}
Speech itself is a complex, continuous valued analogue signal and the most common and natural way to communicate and transmit information between human beings.
Therefore, to enable direct interaction with machines and other non-human systems via speech signals can be highly beneficial for accessability and ubiquituous applications that adress a wide user range.

On the other hand the nature of speech signals also poses difficulties, because conversion and processing in digitalized form is non-trivial and requires multiple (pre-)processing steps.
Typically the majority of these steps are performed using either deterministic mathematical models or neural networks, that are learned based on some kind of training data in a supervised manner.
This causes another issue, because the process of producing and labeling high quality datasets is time- and cost-intensive.
Therefore, it is usually impossible to create and train a fully optimized model for every new task.
Exactly this problem arises for example with relatively unknown or not heavily used languages. 

Rather than working through all development steps, often already existing more general systems are used and fine-tuned for a specific specialized task.
This is also the case for this project, because the provided dataset contains only 6000 training samples with around 10.5 hours of speech, which is not sufficient to train a high quality system with sufficient performance.
Additionally, the available resources especially in case of processing capacity and hardware are very limited, which leads to the next problem. 
Nowadays speech recognition models typically have between several hundred million to a few billion parameters and require datsets of 1000 to 10'000 hours of transcribed speech and weeks of training runs on dedicated hardware to achieve good results.
On the other hand 10 to 100 hours can be enough to adapt respectively fine-tune existing options to a new language or other tasks, because most of the needed feature and mapping functionalities can be already generalized properly \cite{choudhary2026speech}.

The advantage of using Esperanto in this case is the history of the language itself.
As mentioned before it was artificially made and therefore heavily relies on known patterns.
Many words are similar to or derived from western european languages like french, english, german and latin.
Other parts are heavily influenced by slawic langugaes like russian and polnish, because of the polish roots of the inventor.
A big part of the phonemes and their according transcriptions can because of this be assumed to be quiet similar to forementioned languages, for which a lot of training data is already publicly available.
Most commercial and popular open-source models are trained with a lot of data from these domains and should therefore be already able to transcribe the basics of Esperanto. 

Additionally, the grammatical and orthographical rules of Esperanto are by purpose kept very simple.
Back then one major design goal was to make it easy to learn, what resulted in a language which is phonematic, meaning that it has a bijective mapping between phonemes and symbols.
To have a sufficient variety of phonemes avilable these symbols consitst out of latin symbols but are appended with additional diacritic extensions.
Particulary these special symbols could pose a problem to the models because they were not part of the original training data. 


Furthermore, speech recognition itself is a demanding task not only because of the variety of languages, but also because of the required quality of the recorded audio.
Normal usage scenarios for speech recognition systems barely include a perfectly controlled, quiet environment with high quality recording but rather different speakers, a wide variety of background noises and low-quality recording devices.
So a fully functioning speech recognition system must also include preprocessing steps to minimize such effects. 
This can either be done deterministically with classical signal processing techniques or once again with end-to-end systems based on neural networks.
Achieving a sufficient performance in arbitrary environments therefore requires even more data samples to cover as much of these additional effects as possible.

So in total the design of a speech recognition system is a task that includes a lot of steps and possible difficulties that have to be encountered.
Some possible solutions and their tradeoffs will be discussed in the following \autoref{solution}.